\documentclass[11pt,a4paper]{article}
\usepackage[margin=23mm]{geometry}
\usepackage{amsmath,amssymb,booktabs,array,longtable,hyperref}
\hypersetup{hidelinks}
\emergencystretch=3em
\newcommand{\Lc}{\textbf{(c)}\ }
\newcommand{\Ld}{\textbf{(d)}\ }
\newcommand{\dd}{\mathrm d}
\newcommand{\CC}{\texttt{onebody\_\allowbreak continuum.py}}
\newcommand{\TT}{\texttt{onebody\_\allowbreak toda.py}}
\title{A one-body Toda--Boussinesq dispersion test:\\
the bad-equation continuation, its resonance, and a numerical mismatch}
\author{Research note for Gustav Delius}
\date{4 November 2026}
\begin{document}
\maketitle
\begin{abstract}
\Lc\ With the inherited classical units and exact dressed momentum, the
saved Toda data give a fitted small-amplitude one-loop coefficient
$F_T/|P_B|=0.15912\pm0.00334$. The uncertainty is a deliberately broad
fit-sensitivity envelope, not a rigorous error bound. Direct group-velocity
data give $0.15906\pm0.00997$. The good-equation benchmark is $-2/\pi$.
A dimensionally consistent \emph{formal continuation} gives instead
$+2/\pi=0.63661977$ for the bad equation: a discrepancy
$0.47750\pm0.00334$, approximately a factor of four.

\Lc\ The real stable-band trace of the bad equation is different again:
it adds $-\frac{\sqrt3}{2}(\sinh^2\beta+\frac34)$ to the formally continued
answer. Its normal-ordering energy has sign $-3\sinh\beta/\pi$, correcting
the opposite sign in the older bad-equation note. A differential combination
removes constant energy, momentum-origin and mass-counterterm ambiguities
and retains the same cubic coefficient $+2/\pi$.

\Ld\ There is no bad-equation real-frequency trace over the full continuum
with just the good-equation assumptions. Neither continuation nor deletion
of unstable modes follows from those assumptions. Thus the results disfavor
the specified conditional targets within the tested numerical extrapolations;
they do not refute a completed quantum continuum theory or distinguish the
mode and canonical charge pairs. The physical charge identification remains
inconclusive because the determinant/domain and quantum continuum matching
are not established, and the required interchange of limits is unproved.
\end{abstract}

\section{What is being tested, and what is conditional}
Labels are (a) literature proof, (b) literature heuristic, (c) project
derivation/check, and (d) unestablished identification. All new mathematical
formulas below are checked by \CC\ or \TT, as indicated by section.
Exact algebra and numerical evidence are distinguished. The scripts use
SymPy and 65-digit mpmath; high precision in the data analysis does not
increase the precision of the inherited double-precision sea solves.

The inputs read are the complete internal notes
\texttt{det\_\allowbreak projected\_\allowbreak oneloop.tex},
\texttt{frame\_\allowbreak dictionary.tex},
\texttt{toda\_route.tex}, and \texttt{oneloop.tex}, together with the
relevant checking implementations and saved \texttt{frame\_*.json} data.
No external full text or abstract was read for this note; no external
reference, DOI, arXiv identifier, or literature theorem is asserted.
In particular the old note's identification of a stable-band quantum domain
with an inverse-scattering reduction is not adopted as a theorem.

\paragraph{Loop order and fixed momentum. \Lc\ (\CC, \TT).}
We write
\begin{equation}
 E(P)=E_0(P)+\hbar F(P)+O(\hbar^2),\qquad
 F=e-vp=-L_1,\qquad E_0'=vP_0'. \label{eq:invariant}
\end{equation}
Thus $F$ is the order-$\hbar^0$ contribution to $E/\hbar$, and the
order-$\hbar$ contribution to energy. A tangent relabeling
$(e,p)\mapsto(e+v\delta p,p+\delta p)$ leaves $F$ invariant.
For the canonical pair $p_{\rm can}=\partial_v L_1$,
$e_{\rm can}=vp_{\rm can}-L_1$, the same $F$ follows.
Consequently a one-body fixed-momentum experiment cannot decide which of
these two pairs represents the renormalized charges.

\paragraph{Exact qualifications inherited from the det package.}
The good benchmark uses its specified vacuum Wick ordering, no added finite
multiple of $M$, and an equal-mode-count trace \emph{including} the cutoff
endpoint. Dropping that endpoint leaves an uncanceled linear divergence.
\Ld\ The identification of this trace with the absolutely normalized
projected determinant is not proved. The instantaneous symplectic projector
does not preserve fixed potential endpoints; shift and width tangents have
divergent kinetic norms. A common boundary/constraint evolution problem,
its Schur complement, nonzero-frequency completeness, and time-endpoint
determinant phase remain missing. Differentiating the on-shell $L_1$ changes
the endpoint sector and does not prove the canonical charge identification.
The good real-energy benchmark is restricted to the light branch
$0<\alpha<\pi/3$. Nothing in this note resolves these issues.

\paragraph{No new physical prescription is selected.}
The older \texttt{oneloop.tex} records a PI choice of stable-band reduction
and an additional momentum-endpoint prescription. The present request asks
for the same method as the newer good-equation package. These are distinct
inputs. Below we evaluate both the formal continuation and a fully stated
stable-band diagnostic; neither is silently selected as the physical bad
determinant. An additional prescription would require the PI's choice.

\section{Bad versus good: the dictionary and the sign}
\paragraph{Classical dictionary. \Lc\ (\CC).}
For the bad equation the stipulated symplectic normalization gives
\begin{align}
 \mathcal L_b&=\tfrac12\phi_t^2-\tfrac12\phi_x^2-\phi_x^3
                         +\tfrac12\phi_{xx}^2, &
 u_0&=\tfrac{k^2}{2}\operatorname{sech}^2(k\xi/2),\\
 k&=\sinh\beta, &v&=\epsilon\cosh\beta,\qquad
 \ell_b=\epsilon\cosh3\beta,\\
 M&=2k, &\Pi&=-2vk,\qquad P_0=\tfrac23vk^3,\\
 E_0&=\tfrac23k^3+\tfrac8{15}k^5, &L_0&=vP_0-E_0=\tfrac2{15}k^5.
 \label{eq:badclass}
\end{align}
The good values are $k_g=\sin\alpha$, $v_g=\epsilon\cos\alpha$,
$\ell_g=-\epsilon\cos3\alpha$, $E_{g,0}=2k_g^3/3-8k_g^5/15$.
Formally $\alpha=i\beta$ gives $k_g=ik$, $v_g=v$, and
$\ell_g=-\ell_b$. The sign of the spectral parameter is part of the
dictionary; it cannot be omitted.

\paragraph{Why substituting $\alpha=i\beta$ alone is insufficient.
\Lc\ (\CC).}
Introduce the algebraic family
\begin{equation}
 \mathcal L_\sigma=\tfrac12\phi_t^2-\tfrac12\phi_x^2
             +\sigma\phi_x^3-\tfrac\sigma2\phi_{xx}^2,
 \qquad \sigma=c^2.
\end{equation}
For positive $c$, put $(x,t)_\sigma=c(x,t)_g$ and
$\phi_\sigma=\phi_g/c$. Direct substitution gives
\begin{equation}
 S_\sigma=c^{-2}S_g,\qquad \hbar_g=c^2\hbar_\sigma,
 \qquad (E_0,P_0)_\sigma=c^{-3}(E_0,P_0)_g,
 \qquad F_\sigma=c^{-1}F_g. \label{eq:scale}
\end{equation}
Continue these algebraic relations to $c=i$, using branches continuous
from the light good soliton. In particular $\hbar_g=-\hbar_b$ and
$u_g=-u_b$ in this \emph{formal} calculation. This is not a real negative
Planck constant or a unitary equivalence of quantum theories.
\Ld\ Continuing the renormalized trace without extra contour contributions
is an additional assumption, not a consequence of the classical dictionary.

\paragraph{The formal continued target.}
\Lc\ Conditional on that assumption (\CC).
The analytic good fixed-$P$ expression on its light branch is
\begin{equation}
 F_g=\frac{\sqrt3}{\pi}\left(\frac34-k_g^2\right)
        \operatorname{artanh}\frac{2k_g}{\sqrt3}
           -\frac{3k_g}{2\pi}.
\end{equation}
With \eqref{eq:scale} its bad analogue is
\begin{equation}
 \boxed{F_{\rm ac}(\beta)=-iF_g(i\beta)
   =\frac{\sqrt3}{\pi}\left(k^2+\frac34\right)
          \arctan\frac{2k}{\sqrt3}-\frac{3k}{2\pi}.} \label{eq:Fac}
\end{equation}
It obeys $L_{1,\rm ac}=-F_{\rm ac}$ and
\begin{align}
 p_{\rm can,ac}&=-\frac{2\epsilon\sqrt3}{\pi}\cosh\beta
                      \arctan\frac{2\sinh\beta}{\sqrt3},\\
 F_{\rm ac}&=\frac4{3\pi}\beta^3+\frac{14}{45\pi}\beta^5+O(\beta^7),
 &\frac{F_{\rm ac}}{|P_0|}&\longrightarrow+\frac2\pi. \label{eq:aclimit}
\end{align}
The good answer instead has
$F_g=-4\alpha^3/(3\pi)+14\alpha^5/(45\pi)+O(\alpha^7)$ and ratio
$-2/\pi$. Thus even this specific continuation does \emph{not} give the
same small-amplitude coefficient. Forgetting the action/Planck scaling
would reverse the continued sign. Shared long-wave classical behavior
does not by itself identify a trace integrating all wavelengths.

\section{The real bad-equation trace and its extra resonance}
All calculations in this section are \Lc\ (\CC); interpreting their trace
as a projected quantum energy remains \Ld.

\paragraph{The full-line obstruction.}
Linear vacuum modes have $\omega^2=q^2-q^4$. For $|q|>1$ there is no real
positive-frequency oscillator vacuum of the good-package kind.
For example the continued Hirota transmission has modulus
$1.8388478863$ at $\beta=0.7,q=1.3,\epsilon=+1$.
A real scattering phase over the entire line is therefore unavailable.
There is no unique bad analogue \emph{with only} the same conditional
ingredients as the good package.

\paragraph{Explicit stable-band diagnostic.}
Restrict to $|q|<1$ with $\omega=|q|\sqrt{1-q^2}$.
The Hirota polynomial and transmission are
\begin{align}
 \mathcal H(K,W)&=W^2-K^2-K^4,\\
 A(q)&=-\frac{\mathcal H(k-iq,kv-i\omega)}
                    {\mathcal H(k+iq,kv+i\omega)}
       =\frac{Z+iX}{Z-iX},\\
 X&=3kq,\qquad
 Z=v\,\operatorname{sgn}(q)\sqrt{1-q^2}-1-2k^2+2q^2.
 \label{eq:A}
\end{align}
The potential mode $f=\partial_\xi[e^{iq\xi}(1+Ae^{k\xi})/(1+e^{k\xi})]$
satisfies the bad potential pencil
\begin{equation}
 [-\partial_\xi^4+(v^2-1)\partial_\xi^2
       -6\partial_\xi(u_0\partial_\xi)
       +2iv\nu\partial_\xi-\nu^2]f=0,
 \qquad \nu=\omega-vq.
\end{equation}
The script independently checks this residual, $|A|=1$ in the stable band,
and the differentiated phase.

Let $q=\sin t$, $-\pi/2<t<\pi/2$, and choose the continuous logarithm on
each half-band with $\theta(0^\pm)=0$ at fixed $\beta>0$. Then
\begin{equation}
 \theta_t=\frac{2(ZX_t-XZ_t)}{Z^2+X^2}.
\end{equation}
Retaining the equal-mode-count endpoint in the co-moving trace would give
\begin{equation}
 F_{\rm band}=\frac1{4\pi}\int_{-1}^{1}(\omega-vq)\theta'(q)\,\dd q
       -\frac{[(\omega-vq)\theta]_{-1}^{1}}{4\pi}
       -\frac{3k}{\pi}. \label{eq:bandtrace}
\end{equation}
For the stated continuous branches $\theta(1)+\theta(-1)=0$ and the
displayed endpoint vanishes. This is an explicitly checked branch diagnostic,
not a derivation of finite-box endpoint quantization. Changing half-band
branches can change that endpoint and the momentum origin.

\paragraph{Normal-ordering sign.}
Here $\langle\phi_x^2\rangle/\hbar=1/(2\pi)$ on the stable band, but
the bad Hamiltonian contains $+u^3$, whereas the good one contains $-u^3$.
Hence $:u^3:=u^3-3\langle u^2\rangle u$ gives
\begin{equation}
 F_{\rm ct,b}=-\frac{3M}{2\pi}=-\frac{3k}{\pi}.
\end{equation}
The older bad note's $+3k/\pi$ is the good Hamiltonian's sign.
Using that old sign would add $6k/\pi$ to \eqref{eq:bandtrace}; it is a
mass-counterterm difference, not the stipulated bad vacuum Wick ordering.

\paragraph{Evaluation, including the short-wave term.}
Adding the two half-bands and setting $z=\cos t$ reduces the integral
before its counterterm to
\begin{align}
 I&=\frac{3k}{2\pi}\int_0^1 R_k(z)\,\dd z,\\
 R_k(z)&=\frac{2(-4k^2+8z^4-10z^2-1)}
 {16k^4+16k^2z^2+8k^2+16z^4-8z^2+1}.
\end{align}
An antiderivative vanishing at zero is
\begin{equation}
 z-\frac{\sqrt3(4k^2+3)}{12k}
 \left[\arctan\frac{z}{\sqrt3k}
 +\arctan\frac{\sqrt3z(4z^2+8k^2-1)}{3k(4k^2+1)}\right].
\end{equation}
At $z=1$ the sum of angles equals
$\pi-2\arctan(2k/\sqrt3)$; both sides lie in $(0,\pi)$ and have the
same tangent. Thus
\begin{equation}
 \boxed{F_{\rm band}=F_{\rm ac}
           -\frac{\sqrt3}{2}\left(k^2+\frac34\right).} \label{eq:bandclosed}
\end{equation}
The independent 65-digit quadratures at
$\beta=0.01,0.05,0.2,0.7$, for both orientations, agree with this expression
within $2.4\times10^{-66}$; exact symbolic differentiation checks the
antiderivative. This is a trace identity, not a completeness theorem.

The small-amplitude expansion is
\begin{equation}
 F_{\rm band}=-\frac{3\sqrt3}{8}-\frac{\sqrt3}{2}\beta^2
     +\frac4{3\pi}\beta^3-\frac{\sqrt3}{6}\beta^4
     +\frac{14}{45\pi}\beta^5+O(\beta^6).
 \label{eq:bandexp}
\end{equation}
In particular $F_{\rm band}(0.01)=-0.6496052338424$ and there is no finite
$F_{\rm band}/|P_0|$ limit. The constant
$-3\sqrt3/8=-0.6495190528$ is consistent with the short-wave phase winding:
at $q_*=-\sqrt3/2$ for a right mover,
$\omega(q_*)=\sqrt3/4$, $\omega'(q_*)=1$, and
$-\frac12(\omega(q_*)-q_*)=-3\sqrt3/8$.
The algebraic term in \eqref{eq:bandclosed} makes precise why continuing a
closed good trace and integrating the real bad stable band are different
operations. It cannot be omitted as if analytic continuation automatically
preserved the physical mode contour.

\section{Toda extraction at fixed dressed momentum}
This section is \Lc\ (\TT), numerical where fitted. It uses no two-body
phase and no Phase 2i Wiener--Hopf formula.

\paragraph{Units and the exact one-body identities.}
The inherited dictionary, with $h=\hbar_T$, is
\begin{equation}
 \kappa=\sqrt3\beta,\quad h=3b\hbar_B,\quad b=e^{-a/2},\quad
 P_B=\frac{a\overline P}{6\sqrt3 b},\quad
 E_B=\frac{\overline E}{6\sqrt3 b^2}. \label{eq:units}
\end{equation}
It is a leading classical continuum map. In the saved data $a=1$.
The Fredholm identities are inherited checked inputs from
\texttt{frame\_identities.py}; \TT\ checks their unit and chain-rule
consequences and reanalyzes their saved measurements.
The frame package supplies $D=P'(p)$, $V=E'(p)$ and
the edge-subtracted $(\overline P,\overline E)$, with zero at $p=q$.
Its exact momentum inversion is
\begin{align}
 \overline P(q+hu(y))&=P_{T,0}(\kappa)
     =\frac{4b}{a}(\kappa\cosh\kappa-\sinh\kappa),
 &\kappa^2&=yh/b,\\
 E_{T,0}(\kappa)&=2b^2(\sinh\kappa\cosh\kappa-\kappa),
 &\frac{\dd u}{\dd y}&=\frac{2\sinh\kappa}{aD},\\
 \frac{\dd\overline E}{\dd\overline P}&=\frac VD,
 &\frac{\dd E_{T,0}}{\dd P_{T,0}}&=ab\frac{\sinh\kappa}{\kappa}.
\end{align}
The finite-amplitude classical Toda dispersion is subtracted \emph{before}
any continuum extrapolation. The dictionary reproduces the leading cubic
momentum and energy and the fifth-order excess energy above sound;
it is not an exact finite-$\beta$ identification.

Define, exactly as reconstructed from the stored data,
\begin{align}
 C_{E,h}(y)&=\frac{3}{a}\frac{\overline E-E_{T,0}}{hP_{T,0}},\\
 C_{v,h}(y)&=\frac3a\frac{V/D-ab\sinh\kappa/\kappa}{h}.
 \label{eq:coeff}
\end{align}
Then $E_B-E_{B,0}=\hbar_B P_B C_{E,h}$ and the quantum velocity correction
in the fixed inherited units is $\hbar_B C_{v,h}$.
The factor $3/a$ follows from \emph{all three} scales of energy, momentum
and Planck's constant. It is not $1$, $3b$, or $1/(3b)$.
For a smooth edge limit, the chain rule also gives
$C_v=C_E+(2y/3)C_E'$.

\begin{center}\small
\begin{tabular}{rrrrr}
\toprule
$y$ & $C_E(y)$ & edge spread & $C_v(y)$ & edge spread\\\midrule
4 & 0.160932810 & $2.15\,10^{-6}$ & 0.159681133 & $5.68\,10^{-6}$\\
8 & 0.160019933 & $7.36\,10^{-6}$ & 0.159427446 & $2.03\,10^{-5}$\\
16 & 0.159580964 & $2.63\,10^{-5}$ & 0.159289044 & $7.33\,10^{-5}$\\
32 & 0.159360609 & $9.50\,10^{-5}$ & 0.159193402 & $2.59\,10^{-4}$\\
64 & 0.159216903 & $3.35\,10^{-4}$ & 0.158986499 & $9.68\,10^{-4}$\\
\bottomrule
\end{tabular}
\end{center}
These are new independent fits of the saved finite-sea measurements, not
new sea solves. The script checks the input SHA-256 against the frame
analysis and reproduces its energy extrapolations to $10^{-12}$.

\paragraph{Error budget and large-$y$ extraction.}
The raw $h$ values are $.08,.04,.02,.01,.005$, with the frame package's
two refined meshes at $.01,.005$. Central edge values use a quadratic in
the last three $h$ values. The spread compares the last-two linear,
last-four cubic, and last-four $(1,h,h^2,h\log h)$ fits. Measured mesh
differences are multiplied by the largest edge-fit weight norm, $15$.
We then fit $y\ge4,8,16$ through $64$, with each of four models:
\begin{equation}
 C+A/y,\qquad C+A/y+B/y^2,\qquad
 C+A/\sqrt y+B/y,\qquad C+A/y+B\log(y)/y.
\end{equation}
All coefficients are free; $1/(2\pi)$ is never supplied to a fit.
The central model is $C+A/y$ on $4\le y\le64$.
For each fit, edge and mesh diagnostics are propagated with absolute
intercept weights. The final envelope is the maximum over all twelve fits
of the displacement from the central intercept plus that fit's propagated
edge and mesh diagnostics.

\begin{center}\small
\begin{tabular}{lrr}
\toprule Diagnostic & Energy coefficient & Velocity coefficient\\\midrule
Central intercept & 0.159119876 & 0.159062733\\
Central propagated edge spread & 0.000184053 & 0.000525545\\
Central propagated mesh difference & $3.15\,10^{-9}$ & $4.71\,10^{-9}$\\
Largest intercept model/window displacement & 0.000208568 & 0.001051254\\
Full combined envelope & 0.003339590 & 0.009973693\\
\bottomrule
\end{tabular}
\end{center}
The wider envelope retains the poorly conditioned three-point
$(1,y^{-1/2},y^{-1})$ extrapolation; it is not discarded to sharpen the result.
The fitted coefficient is consistent with $1/(2\pi)=0.1591549431$,
but an exact identification with that number is \Ld.

\paragraph{Order of limits: the unbounded systematic.}
At fixed $y$, $\hbar_B/\beta^2=1/y$, so a finite-$y$ table is a quantum
crossover, not a one-loop coefficient. The desired semiclassical coefficient
takes $h\to0$ at fixed nonzero momentum and then $\beta\to0$.
The data instead permit $h\to0$ at fixed $y$ and a subsequent fitted
$y\to\infty$. A simultaneous limit must keep $hy/b\to0$.
\Ld\ Exchanging these operations requires uniform asymptotic control that
the finite data do not supply. No numerical bound for this systematic,
or for untested large-$y$ behaviors, is included in the displayed envelope.
It would be incorrect to interpret the discrepancy as a many-sigma proof.

\Lc\ As a separate check of fixed, unscaled momentum, the saved data give
the Toda-unit $F_T=0.007775546,0.020319171,0.030098075$ at
$\overline P=0.15,0.4,0.6$, respectively. These finite-amplitude values
reproduce the frame package, but are not sufficiently close to zero to
prove the small-amplitude coefficient by themselves.

\section{Comparison and convention tests}
\paragraph{Numerical comparison. \Lc\ (\CC, \TT).}
Within the fixed units and the tested fit families,
\begin{align}
 C_{\rm ac}&=+2/\pi=0.6366197724,\\
 C_T&=0.1591198755\pm0.0033395904,\\
 C_{\rm ac}-C_T&=0.4774998968\pm0.0033395904.
\end{align}
The nominal ratio is $4.00088$. Reusing the good coefficient without a
dictionary would instead give a difference $-0.7957396479$; it does not
repair the comparison. In the specified stable-band trace the ratio
$F/|P|$ does not have a finite limit at all, by \eqref{eq:bandexp}.

\paragraph{Which constants matter? \Lc\ (\CC).}
A vacuum-energy subtraction or a chemical-potential offset per edge
particle is independent of its momentum and drops out of
$\overline E=E(p)-E(q)$. It cannot supply a term proportional to $|P|$.
An order-$\hbar$ change of energy and momentum origins, and a finite mass
counterterm, act on the coefficient on a right-moving branch as
\begin{equation}
 F(\beta)\longmapsto F(\beta)+A-B\cosh\beta+2C\sinh\beta.
 \label{eq:origins}
\end{equation}
Thus $\dd F/\dd P$ removes $A$, but does \emph{not} by itself remove a
momentum-origin change $B$. Our quoted velocity comparison fixes that origin
to the same Fermi edge as the energy. Replacing $P(q)$ by $0$ only in the
momentum label would be an inconsistent reference change.

A stronger origin- and mass-counterterm-independent diagnostic is
\begin{equation}
 \mathcal I[F]=\lim_{\beta\to0^+}
 \frac{(\partial_\beta^3-\partial_\beta)F}
      {(\partial_\beta^3-\partial_\beta)P_0}.
 \label{eq:I}
\end{equation}
The differential operator annihilates $1$, $\cosh\beta$ and
$\sinh\beta$ exactly, while its denominator tends to $4$.
Consequently
\begin{equation}
 \mathcal I[F_{\rm ac}]=\mathcal I[F_{\rm band}]=2/\pi.
\end{equation}
For a differentiable asymptotic dispersion $F_T=C_TP_0+o(\beta^3)$ with
the corresponding differentiated remainder small, $\mathcal I[F_T]=C_T$.
\Ld\ That differentiable expansion is a conditional interpretation of the
Toda fits, not a measurement of a third derivative at zero.
The independent finite-$y$ velocities provide a less derivative-sensitive
check in the fixed edge convention.

\Lc\ The stable-band resonance difference is
$-\frac{\sqrt3}{2}(v^2-\frac14)$. Its second $v$ derivative is
$-\sqrt3$, so it is not an affine energy/momentum-origin shift.
Constants can remove its zeroth-order limit, and a momentum-origin change
can cancel its leading $\beta^2$ piece, but \eqref{eq:I} still retains
$+2/\pi$. Hence none of the changes in \eqref{eq:origins} can change the
cubic discrepancy. Arbitrary chemical potentials coupled to other
nonlinear soliton charges would change the Hamiltonian being compared;
they are additional matching data, not merely the edge energy zero.

\paragraph{Quantum sound speed and time normalization. \Lc/\Ld\ (\TT).}
The saved Toda sound-speed shift in Boussinesq units is
\begin{equation}
 s_{1,T}=0.2387324171,
 \qquad \text{edge spread }1.69\times10^{-8},\quad
 \text{propagated mesh }5.33\times10^{-10}.
\end{equation}
Subtracting the quantum sound line consistently on the Toda side gives
\begin{equation}
 C_T-s_{1,T}=-0.0796125416\pm0.00333961,
\end{equation}
not $\pm2/\pi$. A change of time normalization or of the background
Hamiltonian can add a term $\gamma |P|$ and therefore alter the leading
coefficient. It would require a relative shift of $0.47750$ to reconcile
these particular targets, and is not a vacuum-energy constant or momentum
origin. The frame package explicitly keeps the inherited classical time
scale; no such rescaling is made here. \Ld\ The good det package does not
compute the renormalized continuum vacuum sound speed or its bad quantum
matching. A comparison after matching both sound speeds is therefore not
completed by subtracting Toda's sound shift alone.

\section{What this establishes, checks, and failures recorded}
\paragraph{Verdict.}
\Lc\ This is evidence against identifying either the formally continued
trace or the explicitly evaluated stable-band trace with the Toda one-body
dispersion in the inherited quantum units and tested asymptotic fit class.
It is a one-body mismatch and cannot be blamed on the undetermined
antisymmetric two-body phase frame. It also shows that the proposed
unchanged $-2/\pi$ target is not a valid bad-equation target.

\Ld\ It is \emph{not} a refutation of the good-equation conditional
charge invariant, nor an unconditional refutation of bad-Boussinesq
quantization. The full bad trace requires an extra prescription; the
projected determinant and quantum charge/background matching remain open;
and finite-$y$ extrapolations do not prove the required iterated limit.
The physical comparison is therefore \textbf{inconclusive}, with a large,
explicit conditional numerical discrepancy. Neither success nor failure
of this one-body test could select $p_{\rm mode}$ over
$p_{\rm can}=\partial_v L_1$: their common invariant is exactly the tested
quantity. A successful test would support only that common invariant.

\begin{center}\small
\begin{tabular}{p{.47\linewidth}cp{.33\linewidth}}
\toprule Claim & Label & Check/status\\\midrule
Action, Planck and good/bad dictionaries & (c) & \CC; symbolic scaling\\
Formal bad $F_{\rm ac}$ and $+2/\pi$ & (c)/(d) & \CC; continuation assumption\\
Stable-band trace and tadpole sign & (c) & \CC; exact primitive and quadrature\\
Complete physical bad determinant & (d) & unstable band/domain unresolved\\
Toda fixed-dressed-$P$ coefficients and velocity & (c) & \TT; stored data reanalysis\\
Large-$y$ coefficient $0.15912\pm0.00334$ & (c) & numerical, model-dependent envelope\\
Origin/mass-counterterm invariant & (c) & \CC; exact annihilator\\
Physical charge identification and limit exchange & (d) & not established\\
\bottomrule
\end{tabular}
\end{center}

\paragraph{Failures recorded.}
\begin{enumerate}
\item A scratch derivation initially inverted the Hirota transmission ratio,
giving the wrong sign of the stable-band trace. It was corrected to
$-\mathcal H(\mathrm{difference})/\mathcal H(\mathrm{sum})$ before the
final checker. The corrected expression was then checked against the
potential fluctuation equation and a symbolic antiderivative; no tolerance
was relaxed.
\item The inherited bad-equation note's positive normal-ordering energy
fails the bad Hamiltonian sign check. The corrected sign is negative;
the difference is explicitly $6k/\pi$, not hidden in an energy constant.
\item A full real-frequency bad trace using only the good det assumptions
cannot be constructed: $|q|>1$ has unstable frequencies. The extra
continuation/reduction assumption is retained as an obstruction, not
replaced by an unannounced regulator.
\item The hoped-for equal small-amplitude coefficients fail for the
specified conditional targets. Both the sign change under formal
continuation and the stable-band resonance are retained.
\item The Toda fits do not establish a uniform semiclassical limit or an
exact $1/(2\pi)$ coefficient. Their envelopes are not rigorous confidence
intervals; the unmatched-limit systematic has no assigned numerical bound.
\item The historical scripts were not rerun: some write existing data, and
this task authorizes new \texttt{onebody\_*} files only. No existing result
is newly certified. Both new checkers pass their stated tests; no final
assertion tolerance was weakened.
\item The first LaTeX layout had two overlong text lines. Break opportunities
and a shorter paragraph heading fixed them before the final rendering.
\item An added symbolic endpoint check initially failed structural equality:
SymPy's factorization returned an unevaluated $0\sqrt3$. Exact simplification
reduced it to zero. Neither the identity nor a tolerance was changed.
\end{enumerate}

\paragraph{Reproduction and artifacts.}
Run from \texttt{phase2/scripts/}:
\begin{verbatim}
PYTHONDONTWRITEBYTECODE=1 python3 onebody_continuum.py
PYTHONDONTWRITEBYTECODE=1 python3 onebody_toda.py
\end{verbatim}
Each writes only its own \texttt{onebody\_*\_results.json} and ends with
\texttt{ALL CHECKS PASSED}. Runtime is seconds to tens of seconds on the
task machine. Input hashes and fit details are saved in the Toda output.
Build this document with \texttt{pdflatex} twice from \texttt{phase2/}.
No existing file, README status line, or git state is changed.
\end{document}
